\documentclass[10pt,a4paper]{amsart}
\usepackage{amsmath,amssymb,amsthm,graphicx}
% For commutative diagrams we use Xy-pic
\usepackage[all]{xy}

% We use 2cell for 2-commutative diagrams.
\xyoption{2cell}
% Theorem environments.
%
\theoremstyle{plain}
\newtheorem{theorem}[subsection]{Theorem}
\newtheorem{proposition}[subsection]{Proposition}
\newtheorem{lemma}[subsection]{Lemma}

\theoremstyle{definition}
\newtheorem{definition}[subsection]{Definition}
\newtheorem{example}[subsection]{Example}
\newtheorem{exercise}[subsection]{Exercise}
\newtheorem{situation}[subsection]{Situation}

\theoremstyle{remark}
\newtheorem{remark}[subsection]{Remark}
\newtheorem{remarks}[subsection]{Remarks}

\numberwithin{equation}{subsection}

% Macros
%
\def\lim{\mathop{\mathrm{lim}}\nolimits}
\def\colim{\mathop{\mathrm{colim}}\nolimits}
\def\Spec{\mathop{\mathrm{Spec}}}
\def\Hom{\mathop{\mathrm{Hom}}\nolimits}
\def\Ext{\mathop{\mathrm{Ext}}\nolimits}
\def\SheafHom{\mathop{\mathcal{H}\!\mathit{om}}\nolimits}
\def\SheafExt{\mathop{\mathcal{E}\!\mathit{xt}}\nolimits}
\def\Sch{\mathit{Sch}}
\def\Mor{\mathop{\mathrm{Mor}}\nolimits}
\def\Ob{\mathop{\mathrm{Ob}}\nolimits}
\def\Sh{\mathop{\mathit{Sh}}\nolimits}
\def\NL{\mathop{N\!L}\nolimits}
\def\CH{\mathop{\mathrm{CH}}\nolimits}
\def\proetale{{pro\text{-}\acute{e}tale}}
\def\etale{{\acute{e}tale}}
\def\QCoh{\mathit{QCoh}}
\def\Ker{\mathop{\mathrm{Ker}}}
\def\Im{\mathop{\mathrm{Im}}}
\def\Coker{\mathop{\mathrm{Coker}}}
\def\Coim{\mathop{\mathrm{Coim}}}

% Boxtimes
%
\DeclareMathSymbol{\boxtimes}{\mathbin}{AMSa}{"02}

%
% Macros for moduli stacks/spaces
%
\def\QCohstack{\mathcal{QC}\!\mathit{oh}}
\def\Cohstack{\mathcal{C}\!\mathit{oh}}
\def\Spacesstack{\mathcal{S}\!\mathit{paces}}
\def\Quotfunctor{\mathrm{Quot}}
\def\Hilbfunctor{\mathrm{Hilb}}
\def\Curvesstack{\mathcal{C}\!\mathit{urves}}
\def\Polarizedstack{\mathcal{P}\!\mathit{olarized}}
\def\Complexesstack{\mathcal{C}\!\mathit{omplexes}}
% \Pic is the operator that assigns to X its picard group, usage \Pic(X)
% \Picardstack_{X/B} denotes the Picard stack of X over B
% \Picardfunctor_{X/B} denotes the Picard functor of X over B
\def\Pic{\mathop{\mathrm{Pic}}\nolimits}
\def\Picardstack{\mathcal{P}\!\mathit{ic}}
\def\Picardfunctor{\mathrm{Pic}}
\def\Deformationcategory{\mathcal{D}\!\mathit{ef}}
\usepackage{cleveref}
\crefname{originallemma}{original lemma}{original lemmas}
\Crefname{originallemma}{Original lemma}{Original lemmas}
\makeatletter
\newlabel{derived-lemma-unbounded-right-derived}{{07K7}{source}{}{sourcepointer-1}{}}
\newlabel{derived-lemma-unbounded-right-derived@cref}{{[originallemma][0][]07K7}{[1][0][]source}{}{}{}}
\makeatother
\crefname{originalequation}{original equation}{original equations}
\Crefname{originalequation}{Original equation}{Original equations}
\makeatletter
\newlabel{equation-base-change}{{07HF}{source}{}{sourcepointer-2}{}}
\newlabel{equation-base-change@cref}{{[originalequation][0][]07HF}{[1][0][]source}{}{}{}}
\makeatother
\crefname{originaltheorem}{original theorem}{original theorems}
\Crefname{originaltheorem}{Original theorem}{Original theorems}
\makeatletter
\newlabel{injectives-theorem-K-injective-embedding-grothendieck}{{079P}{source}{}{sourcepointer-3}{}}
\newlabel{injectives-theorem-K-injective-embedding-grothendieck@cref}{{[originaltheorem][0][]079P}{[1][0][]source}{}{}{}}
\makeatother
\crefname{originalreference}{original reference}{original references}
\Crefname{originalreference}{Original reference}{Original references}
\makeatletter
\newlabel{lemma-cohomological-dimension-proper}{{095U}{source}{}{sourcepointer-4}{}}
\newlabel{lemma-cohomological-dimension-proper@cref}{{[originalreference][0][]095U}{[1][0][]source}{}{}{}}
\makeatother
\crefname{originallemma}{original lemma}{original lemmas}
\Crefname{originallemma}{Original lemma}{Original lemmas}
\makeatletter
\newlabel{morphisms-lemma-morphism-finite-type-bounded-dimension}{{0A3V}{source}{}{sourcepointer-5}{}}
\newlabel{morphisms-lemma-morphism-finite-type-bounded-dimension@cref}{{[originallemma][0][]0A3V}{[1][0][]source}{}{}{}}
\makeatother
\crefname{originalreference}{original reference}{original references}
\Crefname{originalreference}{Original reference}{Original references}
\makeatletter
\newlabel{theorem-proper-base-change}{{095T}{source}{}{sourcepointer-6}{}}
\newlabel{theorem-proper-base-change@cref}{{[originalreference][0][]095T}{[1][0][]source}{}{}{}}
\makeatother
\title{Proper base change for unbounded complexes of Z/nZ-modules}
\author{The Stacks Project Authors}
\date{stacks-project:etale-cohomology.tex:local-snapshot-20261011; GNU FDL 1.2 or later}
\begin{document}
\maketitle
\noindent\textit{Editorial note.} All delivered mathematical text is the author's own, copied from the pinned Stacks Project chapter file etale-cohomology.tex; no statement or proof step was rewritten, reordered or removed.\par\medskip
\noindent\textit{Editorial note.} The counted proof is the author's complete argument: the proof is written out in full: the reduction to a quasi-compact base with the dimension bound, the choice of the K-injective representative, the degreewise vanishing of the higher pushforwards of the base-changed terms, and the resulting identification of complexes.\par\medskip
\noindent\textit{Editorial note.} The label lemma-proper-base-change-mod-n names the counted statement, and the proof environment immediately following it in the source is delivered verbatim as its proof.\par\medskip
\noindent\textit{Editorial note.} The counted statement is carried by the author's own lemma environment in the chapter source.\par\medskip
\noindent\textit{Editorial note.} The statement refers to the displayed base change equation of the chapter and is delivered verbatim; the six cited inputs are kept as external pointers with their canonical Stacks tags.\par\medskip
\noindent\textit{Allowed standard prerequisites (editorial).}\begin{itemize}
\item Finiteness of the cohomological dimension of a proper morphism. Same chapter, the lemma bounding the cohomological dimension of a proper morphism through the fibre dimension, together with Morphisms, the boundedness of the fibre dimension of a morphism of finite type. Supplies the finite cohomological dimension of f\_* and f'\_* in the sense needed for unbounded derived categories.
\item K-injective resolutions with injective terms. Injectives, the theorem providing a K-injective complex of injective modules representing a given complex. Supplies the representative I\textasciicircum{}bullet of E used in the computation.
\item Computation of an unbounded right derived functor on a complex of acyclics, and the usual proper base change theorem. Derived Categories, the lemma allowing the derived functor to be computed on a complex of acyclics, and Same chapter, the proper base change theorem. Supplies the reduction of the computation to f'\_*(g')\textasciicircum{}\{-1\}I\textasciicircum{}bullet and the identification of that complex with g\textasciicircum{}\{-1\}f\_*I\textasciicircum{}bullet.
\end{itemize}\medskip
\ifdefined\hypertarget\hypertarget{sourcepointer-1}{}\fi
\noindent\textit{Source reference.} Original lemma 07K7: sources/stacks-chow/tags.txt, source lines 9647--9647. Prerequisite pointer invoked at proof line(s) 18919; it names a standard fact of a sibling chapter that the counted argument uses, and every step of the counted proof built on it is retained verbatim in the delivered body.\par\medskip
\ifdefined\hypertarget\hypertarget{sourcepointer-2}{}\fi
\noindent\textit{Source reference.} Original equation 07HF: sources/stacks-etale-cohomology/pinned-original-source.tex, source lines 18332--18332. Prerequisite pointer invoked at proof line(s) 18893; it names a standard fact of the same chapter that the counted argument uses, and every step of the counted proof built on it is retained verbatim in the delivered body.\par\medskip
\ifdefined\hypertarget\hypertarget{sourcepointer-3}{}\fi
\noindent\textit{Source reference.} Original theorem 079P: sources/stacks-chow/tags.txt, source lines 9279--9279. Prerequisite pointer invoked at proof line(s) 18925; it names a standard fact of a sibling chapter that the counted argument uses, and every step of the counted proof built on it is retained verbatim in the delivered body.\par\medskip
\ifdefined\hypertarget\hypertarget{sourcepointer-4}{}\fi
\noindent\textit{Source reference.} Original reference 095U: sources/stacks-etale-cohomology/pinned-original-source.tex, source lines 18786--18786. Prerequisite pointer invoked at proof line(s) 18904; it names a standard fact of the same chapter that the counted argument uses, and every step of the counted proof built on it is retained verbatim in the delivered body.\par\medskip
\ifdefined\hypertarget\hypertarget{sourcepointer-5}{}\fi
\noindent\textit{Source reference.} Original lemma 0A3V: sources/stacks-chow/tags.txt, source lines 12753--12753. Prerequisite pointer invoked at proof line(s) 18901; it names a standard fact of a sibling chapter that the counted argument uses, and every step of the counted proof built on it is retained verbatim in the delivered body.\par\medskip
\ifdefined\hypertarget\hypertarget{sourcepointer-6}{}\fi
\noindent\textit{Source reference.} Original reference 095T: sources/stacks-etale-cohomology/pinned-original-source.tex, source lines 18610--18610. Prerequisite pointer invoked at proof line(s) 18928; it names a standard fact of the same chapter that the counted argument uses, and every step of the counted proof built on it is retained verbatim in the delivered body.\par\medskip
\begin{lemma}
\label{lemma-proper-base-change-mod-n}
Let $f : X \to Y$ be a proper morphism of schemes. Let $g : Y' \to Y$ be
a morphism of schemes. Set $X' = Y' \times_Y X$ and denote
$f' : X' \to Y'$ and $g' : X' \to X$ the projections.
Let $n \geq 1$ be an integer.
Let $E \in D(X_\etale, \mathbf{Z}/n\mathbf{Z})$.
Then the base change map (\ref{equation-base-change})
$g^{-1}Rf_*E \to Rf'_*(g')^{-1}E$
is an isomorphism.
\end{lemma}

\begin{proof}
It is enough to prove this when $Y$ and $Y'$ are quasi-compact.
By Morphisms, Lemma
\ref{morphisms-lemma-morphism-finite-type-bounded-dimension}
we see that the dimension of the fibres of
$f : X \to Y$ and $f' : X' \to Y'$ are bounded. Thus
Lemma \ref{lemma-cohomological-dimension-proper} implies that
$$
f_* :
\textit{Mod}(X_\etale, \mathbf{Z}/n\mathbf{Z})
\longrightarrow
\textit{Mod}(Y_\etale, \mathbf{Z}/n\mathbf{Z})
$$
and
$$
f'_* :
\textit{Mod}(X'_\etale, \mathbf{Z}/n\mathbf{Z})
\longrightarrow
\textit{Mod}(Y'_\etale, \mathbf{Z}/n\mathbf{Z})
$$
have finite cohomological dimension in the sense of
Derived Categories, Lemma \ref{derived-lemma-unbounded-right-derived}.
Choose a K-injective complex $\mathcal{I}^\bullet$
of $\mathbf{Z}/n\mathbf{Z}$-modules each of whose
terms $\mathcal{I}^n$ is an injective sheaf of
$\mathbf{Z}/n\mathbf{Z}$-modules representing $E$.
See Injectives, Theorem
\ref{injectives-theorem-K-injective-embedding-grothendieck}.
By the usual proper base change theorem we find
that $R^qf'_*(g')^{-1}\mathcal{I}^n = 0$ for $q > 0$, see
Theorem \ref{theorem-proper-base-change}.
Hence we conclude by 
Derived Categories, Lemma \ref{derived-lemma-unbounded-right-derived}
that we may compute $Rf'_*(g')^{-1}E$ by the complex
$f'_*(g')^{-1}\mathcal{I}^\bullet$. Another application
of the usual proper base change theorem shows that
this is equal to $g^{-1}f_*\mathcal{I}^\bullet$ as desired.
\end{proof}

\end{document}
